\begin{table}[htbp]
\centering
\caption{Robustness Checks}
\label{tab:robustness}
\begin{threeparttable}
\small
\begin{tabular}{lccccc}
\hline\hline
 & Coefficient & SE & $p$-value & Forces & $N$ \\
\hline
\multicolumn{6}{l}{\textit{Panel A: Baseline}} \\
Baseline (Table 2, col 1) & -2.011 & (4.006) & 0.616 & 42 & 462 \\
\hline
\multicolumn{6}{l}{\textit{Panel B: Placebo Outcome}} \\
Burglary rate & 1.157 & (1.017) & 0.255 & 42 & 461 \\
\hline
\multicolumn{6}{l}{\textit{Panel C: Sample Restrictions}} \\
Drop Greater Manchester & -1.160 & (4.487) & 0.796 & 41 & 451 \\
Drop London & -2.009 & (4.007) & 0.616 & 41 & 451 \\
England only (drop Wales) & -4.903 & (4.282) & 0.252 & 38 & 418 \\
Short post-period (2 years) & 1.268 & (3.773) & 0.737 & 42 & 378 \\
\hline
\multicolumn{6}{l}{\textit{Panel D: Alternative Treatment}} \\
Binary (top quartile ASBO) & -11.361 & (11.024) & 0.303 & 42 & 462 \\
\hline
\multicolumn{6}{l}{\textit{Panel E: Inference}} \\
Permutation $p$-value (1,000 draws) & \multicolumn{5}{c}{0.638} \\
\hline\hline
\end{tabular}
\begin{tablenotes}[flushleft]\small
\item \textit{Notes:} All regressions include force area and year-month fixed effects. Standard errors clustered by police force area. The baseline is the specification from Table \ref{tab:main}, column (1). Panel B uses burglary rate as a placebo outcome (the ASB toolkit consolidation should not directly affect burglary enforcement). Panel C tests sensitivity to influential observations and sample composition. Panel D replaces continuous treatment with a binary indicator for top-quartile pre-reform ASBO issuance. Panel E reports a Fisher exact permutation $p$-value from 1,000 random reassignments of ASBO rates across force areas. $^{*}p<0.10$, $^{**}p<0.05$, $^{***}p<0.01$.
\end{tablenotes}
\end{threeparttable}
\end{table}
